\documentclass[11pt,a4paper]{article}
\usepackage[margin=27mm]{geometry}
\usepackage[T1]{fontenc}
\usepackage{lmodern,amsmath,amssymb,amsthm,booktabs,longtable,array,needspace,xurl,hyperref}

% Temporary author-review layer requested by Jordi. Omitted in main.pdf.
\usepackage{xcolor}
\usepackage[protrusion=true,expansion=false]{microtype}
\newif\ifjordireview
\ifdefined\JordiReview\jordireviewtrue\else\jordireviewfalse\fi
\newcounter{jordinote}
\newcommand{\jordinote}[2]{%
  \ifjordireview
  \stepcounter{jordinote}\par\medskip\noindent
  \fcolorbox{blue!55!black}{blue!4}{%
    \begin{minipage}{\dimexpr\linewidth-2\fboxsep-2\fboxrule\relax}
    \small\raggedright\color{blue!45!black}
    \textbf{Solo Jordi -- J\thejordinote: #1}\par\smallskip
    #2
    \end{minipage}}\par\medskip
  \fi
}

\hypersetup{colorlinks=true,linkcolor=blue!35!black,citecolor=blue!35!black,urlcolor=blue!35!black}
\setlength{\emergencystretch}{3em}
\newtheorem{lemma}{Lemma}[section]
\newtheorem{proposition}[lemma]{Proposition}
\newtheorem{assumption}[lemma]{Assumption}
\newcommand{\bits}{\{0,1\}}
\newcommand{\Enc}{\mathsf{AES}}
\newcommand{\Upd}{\ensuremath{\mathsf{Update}}}
\newcommand{\DF}{\ensuremath{\mathsf{Block\_Cipher\_df}}}
\newcommand{\fall}[2]{(#1)_{#2}}
\newcommand{\hm}{H_{\min}}
\newcommand{\pad}{\operatorname{pad}}
\numberwithin{equation}{section}
\clubpenalty=10000
\widowpenalty=10000
\displaywidowpenalty=10000
\newcolumntype{L}[1]{>{\raggedright\arraybackslash}p{#1}}
\theoremstyle{definition}
\newtheorem{definition}{Definition}[section]
\theoremstyle{plain}

\title{CTR\_DRBG with AES-256:\\An Algorithmic Conformity Analysis for DRG.3}
\author{}
\date{}
\makeatletter
\renewcommand{\maketitle}{%
  \begin{center}{\LARGE\@title\par}\end{center}
  \vspace{.5\baselineskip}}
\makeatother

\begin{document}
\maketitle
\ifjordireview
\begin{center}\small\color{blue!45!black}
\textbf{VERSIÓN DE REVISIÓN PARA JORDI -- COMENTARIOS DE TRABAJO}\\
Los recuadros azules no forman parte del artículo.
\end{center}
\fi
\begin{abstract}
We develop an algorithmic conformity argument for CTR\_DRBG with AES-256 under the DRG.3 requirements of AIS 20/31. The construction is described by its request and seeding tuples, and each required property is related to explicit statements about cipher evaluations and derived states. For the configuration without a derivation function and with fresh conditionally uniform 384-bit seed input, we derive classical prediction and inversion bounds in a shared-key permutation model over up to $2^{48}$ requests. The argument admits public adaptive additional input and accounts for dependence between generated output and the next key. It also establishes effective-state accounting, the entropy effect of initialization and a statistical-distance bound. Each result is linked to its DRG.3 requirement. Source admissibility is assumed. The model-based conclusions are accompanied by a cryptographic assessment and its explicit limits for concrete AES.
\end{abstract}

\jordinote{Qué queremos demostrar}{Jordi: revisa que el objetivo sea el que buscamos: un argumento de conformidad algorítmica DRG.3, con fuente suficiente asumida, siguiendo el recorrido expositivo del ejemplo AIS de Hash\_DRBG. El abstract debe prometer exactamente lo que luego justificamos. Las hipótesis sobre AES y las obligaciones pendientes no deben confundirse con requisitos ya cerrados.}

\noindent\textbf{Keywords:} CTR\_DRBG; AIS 20/31; DRG.3; algorithmic conformity; state compromise.

\section{CTR\_DRBG with AES-256}\label{sec:algorithm}
CTR\_DRBG is the block-cipher-based deterministic random bit generator specified in SP 800-90A Rev.1 \cite{90a}. Its persistent state consists of a cipher key $K$, a counter $V$ and a public reseed counter $c$. A request produces counter-mode cipher outputs and then replaces both $K$ and $V$. The conformity question therefore concerns not only the unpredictability of a counter-mode stream, but also the relation between that stream and the state derived from its subsequent cipher evaluations.

We consider AES-256, with a 128-bit block and counter and a seed length of 384 bits. Counter increments use the full 128-bit register. The derivation function is disabled. A successful request returns $1\le p\le2^{19}$ bits, and at most $R=2^{48}$ requests are permitted between reseeds. Personalization and additional input contain at most 384 bits. These are the fixed parameters of the construction analyzed below.

The operation Update evaluates three successive counter values under the current key, concatenates the resulting 384 bits, XORs the provided data and splits the result into the next key and counter. Instantiate applies Update to the zero state after combining entropy input and personalization. Reseed applies the corresponding operation to the current state and resets $c$. Generate optionally applies Update before producing output, emits the required counter-mode blocks, discards any unused suffix, and always applies Update after output generation. The same processed additional input is used in both updates. Section~\ref{sec:tuples} gives the exact describing maps.

AIS 20/31 supplies a conformity argument for Hash\_DRBG \cite{ais}, section 5.3.1, paragraphs 1009--1076, based on describing tuples, justified cryptographic hypotheses and a separate argument for each required property. Müller \cite{browser}, sections 3.2.4 and 4.2.4, gives the corresponding CTR\_DRBG tuples and a plausibility argument. We use that formulation to develop the algorithmic implications explicitly. The distinction between output prediction, state inversion and the instant of compromise is also central to the analysis of Woodage and Shumow \cite{ws}.

\section{DRG.3 requirements and proof obligations}\label{sec:properties}
The DRG.3 requirements concern three different matters: the source supplying seed material, the amount of effective state and initialized entropy, and the algorithmic behavior of generated output. Table~\ref{tab:obligations} states the obligations considered here. Source admissibility is a premise of the algorithmic analysis; its implementation-specific verification is separate.

\begingroup\small\setlength{\tabcolsep}{5pt}
\begin{longtable}{@{}L{17mm}L{126mm}@{}}
\caption{DRG.3 obligations for the selected source route and algorithmic profile.}\label{tab:obligations}\\
\toprule DRG.3 & Obligation\\\midrule\endfirsthead
\toprule DRG.3 & Obligation\\\midrule\endhead
\bottomrule\endfoot
DRG.3.1 & Seed material is supplied by an admissible source. We select the TRNG route and assume source admissibility.\\[2pt]
DRG.3.2 & A request returns at most $2^{19}$ bits.\\[2pt]
DRG.3.3 & The effective internal state has at least 246 bits. Stored components recoverable once a key candidate is fixed add no independent search factor.\\[2pt]
DRG.3.4 & The initialized effective state has at least 240 bits of min-entropy, or at least 250 bits of Shannon entropy with the required stationary raw output.\\[2pt]
DRG.3.5 & Observed internal random numbers do not enable practically useful prediction of their successor.\\[2pt]
DRG.3.6 & Observed internal random numbers do not enable practically useful reconstruction of their predecessor.\\[2pt]
DRG.3.7 & Disclosure of the current state does not make earlier output practically recoverable at the specified compromise boundary.\\[2pt]
DRG.3.8 & Controlled additional input preserves the required strength.\\[2pt]
DRG.3.9 & Output and state-transition functions are cryptographic; the complete state transition is one-way.\\[2pt]
DRG.3.10 & Appropriate output-only statistical tests behave as for ideal random output.\\
\end{longtable}
\endgroup

Let $Z_1,Z_2,\ldots\in\bits^{128}$ denote internal random numbers in chronological order, with $Z_{y,j}$ denoting block $j$ of request $y$, $Y_y$ the returned bits of request $y$, and $s_{y+1}$ the persistent state after its final Update. A request concatenates blocks and may truncate the last one. Let
\[
 \mathcal V_{\ell:u}=(Z_\ell,\ldots,Z_u,\mathbf a,\mathbf p)
\]
denote the observed blocks together with admitted public inputs, lengths and configuration. A predictor of $n$ distinct withheld bits has baseline success $2^{-n}$ for independent uniform output. Each target bit is withheld before observation; adaptive target selection and public controls may inspect only disclosed information.

\begin{definition}[Forward secrecy, DRG.3.5]\label{def:forward}
Given $\mathcal V_{\ell:u}$ and no secret state, predicting the successor $Z_{u+1}$ must be practically infeasible beyond the guessing baseline. This applies to individual internal numbers, including across a request boundary.
\end{definition}
\begin{definition}[Backward secrecy, DRG.3.6]\label{def:backward}
For $\ell>1$, given $\mathcal V_{\ell:u}$ and no secret state, reconstructing or predicting the predecessor $Z_{\ell-1}$ must be practically infeasible beyond the guessing baseline.
\end{definition}
\begin{definition}[Enhanced backward secrecy, DRG.3.7]\label{def:enhanced}
Given the current state $s_{y+1}$, the actual last additional input and public observations, withheld bits of $Y_1,\ldots,Y_y$ must remain practically unrecoverable. We select request-level granularity, with
\[
 \mathrm{Update}_y\prec\mathrm{Use}(Y_y),
 \qquad \mathrm{erase}(\text{recoverable old and temporary state}).
\]
These are conceptual atomicity and erasure as in AIS paragraphs 109--112 and 134. Disclosure during an unfinished request is a different task.
\end{definition}
AIS paragraph 214 sets an approximately $\min\{240,n\}$-bit classical work benchmark for $n$ old bits and also requires quantum attacks to remain infeasible. The mathematical bounds below address the classical part. Knowledge of the current state and future public inputs determines future output of a deterministic generator; that is not excluded by ordinary forward secrecy.

\begin{definition}[Controlled additional input, DRG.3.8]\label{def:additional}
Let $\mathcal V_y$ be the disclosed history before request $y$. For every admitted policy $a_y=\mathcal A_y(\mathcal V_y)$ choosing its input before use, the strength of the generator, including secrecy and transition one-wayness, must be preserved. No entropy or confidentiality of $a_y$ is assumed. In particular, we claim no entropy refresh through Generate additional input. That separate use as cryptographic post-processing for a PTG.3 source, discussed for Hash\_DRBG in AIS paragraph 1076, is outside the present DRG.3 argument.
\end{definition}
\begin{definition}[Cryptographic functions and one-way transition, DRG.3.9]\label{def:oneway}
The output and transition functions must be built from cryptographic primitives; XOR, projection or increment alone is insufficient. Their suitability still requires analysis. For the actual known parameters $a,p$, inversion means finding any member of
\[
 \phi_{a,p}^{-1}(s^+)=\{s:\phi(s,a,p)=s^+\}.
\]
Finding any such member must be computationally impracticable. The input may not be replaced after the target is seen; the inverse need not be the original predecessor. This is the state-compromise interpretation of AIS paragraph 212.
\end{definition}
\begin{definition}[Statistical inconspicuousness, DRG.3.10]\label{def:statistics}
Tests without secret-state knowledge should behave as on ideal output. We use statistical distance as sufficient model-based evidence:
\[
 |\Pr[T(Y)=1]-\Pr[T(U)=1]|\le\Delta(Y,U),\qquad
 \Delta(P,Q)=\tfrac12\sum_x|P(x)-Q(x)|,
\]
where $U$ is ideal output under the same public request rule.
\end{definition}

\jordinote{Qué sabe el observador}{Antes de entrar en las tuplas, contrasta estas definiciones con AIS 207--215. La orientación temporal y el instante de compromiso importan: forward secrecy ordinaria no promete ocultar el futuro a quien conoce el estado actual. Comprueba también que estamos protegiendo números internos, y no solamente peticiones completas, salvo la granularidad expresamente elegida para .7.}

\Needspace{9\baselineskip}
\section{Describing tuples}\label{sec:tuples}
\Needspace{5\baselineskip}
\subsection{The request 9-tuple}
Identify 128-bit strings with $\mathbb Z_N$, $N=2^{128}$; counter additions are modulo $N$. Write $\Vert$ for concatenation, $\oplus$ for bitwise XOR, and $\operatorname{left}_p$ for the first $p$ bits of a string. The operation $\operatorname{split}_{256,128}$ separates a 384-bit string into its 256-bit prefix and 128-bit suffix. For $x\in\bits^{384}$ define
\begin{align}
 B_i(K,V)&=\Enc_K(V+i),\qquad i\ge1,\label{eq:block}\\
 T(K,V)&=B_1(K,V)\Vert B_2(K,V)\Vert B_3(K,V),\label{eq:triple}\\
 \Upd(x,K,V)&=\operatorname{split}_{256,128}\bigl(T(K,V)\oplus x\bigr).\label{eq:update}
\end{align}
These are identities, not modeling assumptions. AES-256 has seedlen 384; three increments and encryptions supply exactly the string used by the NIST Update algorithm (SP 800-90A, section 10.2.1.2).

Following the describing formulation in \cite{browser}, sections 3.2.4.1 and 4.2.4.1, we take $S_{\rm req}=S$, with auxiliary functions $f,g$ and the decomposition $\phi=\phi_B\circ\phi_A$ at fixed public parameters. All algorithmic parameters are fixed as in Section~\ref{sec:algorithm}. The spaces are
\begin{align*}
 S&=\bits^{256}\times\bits^{128}\times\{1,\ldots,R+1\},\qquad S_{\rm req}=S,\\
 \mathcal R&=\bits^{128},\qquad
 A=\bigcup_{j=0}^{384}\bits^j,\qquad I=\{1,\ldots,2^{19}\}.
\end{align*}
Write $s=(K,V,c)$, where $c$ is the public reseed counter. Let $\pad$ append zero bits to length 384 and write $a=\varnothing$ for Null. Define $g:A\to\bits^{384}$ and $f:A\times\bits^{256}\times\bits^{128}\to\bits^{256}\times\bits^{128}$ by
\begin{align}
 g(a)&=\begin{cases}0^{384},&a=\varnothing,\\\pad(a),&a\ne\varnothing,\end{cases}\label{eq:g}\\
 f(a,K,V)&=\begin{cases}(K,V),&a=\varnothing,\\
 \Upd(g(a),K,V),&a\ne\varnothing.\end{cases}\label{eq:f}
\end{align}
Define the working state
\begin{equation}
 (\bar K,\bar V)=f(a,K,V).\label{eq:pre}
\end{equation} In particular, a nonempty all-zero input invokes Update, whereas Null does not.

For a successful request let $q=\lceil p/128\rceil$. Following the BSI decomposition, define $\phi_A,\phi_B:S\times A\times I\to S$ on successful states by
\begin{align}
 \phi_A((K,V,c),a,p)&=(f(a,K,V),c),\label{eq:phiA}\\
 \phi_B((K,V,c),a,p)&=(\Upd(g(a),K,V+q),c+1),\label{eq:phiB}\\
 \phi(s,a,p)&=\phi_B(\phi_A(s,a,p),a,p).\label{eq:phi}
\end{align}
Pairs returned by $f$ or Update are flattened when forming a state triple. Thus $\phi_A$ performs the optional pre-update and leaves $c$ unchanged; $\phi_B$ advances the working counter by the number of output blocks, performs the final Update and increments $c$. Its argument is the state after $\phi_A$, not the temporary state after the output loop. The remaining maps are
\begin{align}
 \phi_{\rm req}((K,V,c),a)&=(\bar K,\bar V+1,c),\label{eq:phireq}\\
 \phi_0((k,w,c),a)&=(k,w+1,c),\qquad \psi(k,w,c)=\Enc_k(w).\label{eq:psi}
\end{align}
Equivalently, $\phi_{\rm req}((K,V,c),a)=\phi_0((f(a,K,V),c),a)$. Its definition has no request-length argument. The AIS request 9-tuple is
\[
 (S,S_{\rm req},\mathcal R,A,I,\phi,\phi_{\rm req},\phi_0,\psi).
\]
Here $\phi_{\rm req}:S\times A\to S_{\rm req}$, $\phi_0:S_{\rm req}\times A\to S_{\rm req}$ and $\psi:S_{\rm req}\to\mathcal R$. The loop emits $\psi$ before applying $\phi_0$; its first output uses $\bar V+1$. The public counter remains in the temporary state, although $\phi_0$ and $\psi$ do not use its value. The complete transition is
\begin{align}
 \phi((K,V,c),a,p)&=(K^+,V^+,c+1),\label{eq:phi-result}\\
 (K^+,V^+)&=\Upd(g(a),\bar K,\bar V+q).\label{eq:post}
\end{align}
The same $g(a)$ is used before and after output. Absence suppresses only the pre-update. The complete transition depends on $p$, unlike Hash\_DRBG's transition.

The guard rejects $c=R+1$ before generating output. For a total describing map on $S$, use $\min(c+1,R+1)$ in $\phi_B$ and (\ref{eq:phi-result}); this extension agrees with all successful requests. Invalid inputs are rejected before invoking the maps.

The externally returned bit string is
\begin{equation}
 Y=\operatorname{left}_p\left(\big\Vert_{j=1}^{q}\Enc_{\bar K}(\bar V+j)\right).\label{eq:output}
\end{equation}
The suffix of the last block is discarded rather than retained for another request. Internal random numbers are the 128-bit blocks. The distinction between the temporary increment $\phi_0$ and the full transition $\phi$ will be important for one-wayness.

\jordinote{La tupla debe ejecutar la 90A}{Lee una petición a mano: entrada no vacía activa pre-Update; después salen los bloques; siempre hay post-Update con la misma máscara. Distingue $\phi_0$ de $\phi$: incrementar el contador temporal no es la transición cuyo carácter unidireccional queremos justificar. Comprueba el primer incremento, el índice del último bloque, el truncamiento y la diferencia entre Null y una cadena no vacía de ceros.}

\Needspace{5\baselineskip}
\subsection{The two seeding 4-tuples}
For the configuration without a derivation function, let $SM=\bits^{384}$ and $PS=A$. Instantiate is described by $(SM,PS,S,\phi_{\rm seed})$, where
\begin{equation}
 \phi_{\rm seed}(e,z)=\bigl(\Upd(e\oplus \pad(z),0^{256},0^{128}),1\bigr).\label{eq:seed}
\end{equation}
Flatten the pair returned by Update into the first two components of $S$ and set the third component to 1. The reseeding tuple is $(SM,PS,S,\phi_{\rm reseed})$ with
\begin{equation}
 \phi_{\rm reseed}((K,V,c),e,z)=
 \bigl(\Upd(e\oplus \pad(z),K,V),1\bigr).\label{eq:reseed}
\end{equation}
Here $\phi_{\rm seed}:SM\times PS\to S$ and $\phi_{\rm reseed}:S\times SM\times PS\to S$. The old state is the first argument of the reseeding map, as in the AIS reseeding description. No nonce is used in the no-DF instantiation. Personalization and reseed auxiliary data are public; neither is counted as seed entropy. The source premise must hold even when the old state is known.
The tuple directly follows the NIST procedures: Update is specified in section 10.2.1.2, no-DF instantiation in 10.2.1.3.1, reseeding in 10.2.1.4.1 and generation in 10.2.1.5. The optional first update is $\phi_A$; counter-mode generation is the $\phi_{\rm req},\phi_0,\psi$ loop; the mandatory final update is included in $\phi_B$.

For the probability calculations we use only the cryptographic pair
\begin{equation}
 \pi(K,V,c)=(K,V).\label{eq:projection}
\end{equation}
The counter $c$ is determined by the public request history. Projecting it away in a calculation neither changes an emitted block nor removes secret information from a state-compromise observation.

\jordinote{Qué asumimos en el seeding}{Las 4-tuplas separan la fuente del procesamiento algorítmico. Aquí no intentamos certificar una fuente concreta: suponemos entrada fresca uniforme de 384 bits, incluso condicionada al estado anterior. Revisa que la personalización pública no reciba crédito de entropía y que el argumento posterior de traducción XOR se aplique tanto a Instantiate como a Reseed.}

\jordinote{Cómo usar el puente con la 90A}{Estas correspondencias son la comprobación de fidelidad del modelo: cada mapa debe corresponder a un procedimiento de la 90A y conservar su orden de operaciones. No es una hipótesis criptográfica. Para la revisión, empieza por estas correspondencias y las tuplas; una cota correcta sobre un algoritmo distinto no serviría para nuestro objetivo.}

\jordinote{Por qué usamos la formulación del BSI}{Aquí adoptamos directamente $f$, $g$, $S_{\rm req}=S$ y la descomposición con $\phi_A,\phi_B$ del estudio BSI. La variante sin DF usa el padding de su sección 4.2.4. La proyección $\pi$ aparece solo para calcular con el par secreto; no sustituye la tupla principal. Revisa que cada argumento de plausibilidad que desarrollamos quede ligado a estas mismas funciones y a un requisito concreto.}

\jordinote{El perfil y las fronteras de la tupla}{Revisa las decisiones que fijan el objeto de la prueba: AES-256, ausencia de DF, incremento completo y guardas. El esquema con $f,g$ y $S_{\rm req}=S$ procede del antecedente citado; la convención $R+1$ hace explícito el agotamiento. Cualquier ampliación del perfil debe contrastarse con las hipótesis usadas en los conteos. No basta mantener los mismos símbolos.}

\section{Verification of the DRG.3 properties}\label{sec:verification}
The argument first establishes common estimates for cipher evaluations and derived states. It then applies them to the secrecy requirements, followed by the size, input-control, one-wayness and statistical conditions. Initialization is treated separately, as its source premise must not be inferred from the output analysis.

\Needspace{5\baselineskip}
\subsection{Hypotheses and common estimates}\label{sec:model}
In addition to the fixed profile of Section~\ref{sec:algorithm}, we assume fresh entropy input uniform on $\bits^{384}$ conditional on the public history, the old state and the primitive tables. Public additional input, lengths and stopping decisions may depend adaptively on disclosed information but not on withheld target bits. The bounds apply to classical direct and inverse primitive computations during a finite reseed interval.

For a request $y$ with working state $(\bar K_y,\bar V_y)$ and $q_y$ output blocks,
\begin{align}
 Z_{y,j}&=\Enc_{\bar K_y}(\bar V_y+j),\quad 1\le j\le q_y,\label{eq:request-sequence}\\
 (K_{y+1}\Vert V_{y+1})\oplus g(a_y)
 &=\Enc_{\bar K_y}(\bar V_y+q_y+1)\Vert\Enc_{\bar K_y}(\bar V_y+q_y+2)
 \Vert\Enc_{\bar K_y}(\bar V_y+q_y+3).\label{eq:revealed-triple}
\end{align}
The second identity describes the information revealed by the state after request $y$. Its three cipher values depend on the working key that also generated the request output.

\Needspace{6\baselineskip}
\begin{assumption}[Permutation model for AES-256]\label{ass:aes}
The family $\{\Enc_k\}_{k\in\bits^{256}}$ consists of independent uniform permutations of $\bits^{128}$. All direct and inverse computations under an identical key use the same permutation.
\end{assumption}
This is an idealized model of the primitive, including evaluations under the derived keys. It does not assume the DRBG secrecy properties or make Update a random function: $T(K,V)$ and $T(K,V+1)$ share two blocks. The following derivations retain these dependencies. Applying them to concrete AES-256 additionally requires the cryptographic assessment of AIS paragraphs 97--100 and 215.

\paragraph{Cryptographic interpretation.}\label{sec:primitive-assessment}
We use the model as a cryptographic idealization and assess its application through the attack evidence and access assumptions below. AES is not identified with an independent family of random permutations. Single-key cryptanalysis and related-key cryptanalysis must be distinguished. For example, the full-AES biclique result of Bogdanov, Khovratovich and Rechberger \cite{biclique} gives an AES-256 key-recovery cost of approximately $2^{254.4}$ cipher computations. This historical result is above the 240-bit classical benchmark; it is evidence about a particular attack, not a lower bound on every attack against the construction.

Related-key attacks prevent treating independence across concrete AES keys as unconditional. The attacks of Biryukov and Khovratovich \cite{relatedkey}, section 4.1, use encryption or decryption access under secret keys with attacker-selected relations. In the present interface a public mask changes the key obtained from hidden cipher outputs; it does not supply unrestricted access under a chosen relation to the current secret key. The conditional key law and the shared-table argument below quantify that distinction in the model. They do not prove that every structural property of concrete AES is harmless. The cryptographic hypothesis needed for application is that the counter evaluations and the specified derivation of secret keys expose no exploitable AES structure invalidating these model-based conclusions. This concerns the primitive and its derived-key use; the required DRBG properties are then derived rather than assumed.

For the nonclassical assessment in AIS paragraph 214, generic search over a 256-bit key space has an order-of-magnitude cost of $2^{128}$ Grover iterations before implementation costs; Grassl et al.\ \cite{grassl} analyze the corresponding AES circuits. This is an infeasible generic reference attack, not a universal lower bound for quantum attacks on the derived construction. The classical inequalities below are not extrapolated to coherent primitive queries. Concrete classical conclusions rely on the stated primitive hypothesis. The additional quantum-infeasibility assessment relies on the absence of an applicable structural attack that makes recovery feasible in this secret-derived-key use. A complete quantum composition theorem is not asserted.

\jordinote{Hipótesis sobre AES, no sobre la conclusión}{Modelamos cada clave mediante una permutación ideal independiente, reutilizando la misma tabla cuando se repite la clave. No suponemos que el DRBG ya sea seguro ni que Update sea una función aleatoria independiente: las ternas de contadores vecinos comparten bloques. La revisión debe vigilar que ninguna independencia posterior contradiga estas dependencias.}

Strengthen the observer by giving every full output block except its target bits, public controls and the terminal state. Generation ends at that compromise; at most $Q$ direct or inverse primitive computations may occur before, during or afterwards. Each target remains withheld from the controller.

\paragraph{Conditional distributions.}
For an auxiliary calculation give each group of evaluations a private independent permutation: either three pre-update evaluations, or $q_i$ output evaluations followed by three final-update evaluations. Requests are indexed by $y$ and groups by $i$; $y(i)$ identifies the request containing group $i$. Write
\[
 \sigma_i=(K_i,V_i),\quad W_i=R_i\Vert U_i,\quad |R_i|=q_i,\quad |U_i|=3,
 \quad \sigma_{i+1}=\operatorname{split}_{256,128}(U_i\oplus g(a_{y(i)})).
\]
Both groups of a nonempty-input request therefore use the same actual input $a_{y(i)}$. The pair $\sigma_i=(K_i,V_i)$ is the projection (\ref{eq:projection}); it is not the full state of $S$. Before an output group this pair is $f(a_{y(i)},K,V)$, and its final pair is the projection of $\phi_B$. A pre-update has $q_i=0$. At most $r\le2^{48}$ requests produce $t\le2r$ groups. Let $b=4096$, $L=b+3$, and $\fall{x}{j}=x(x-1)\cdots(x-j+1)$, with $\fall{x}{0}=1$.

The distinct counter inputs imply, for any distinct list $w$ and fixed incoming state,
\begin{equation}
 \Pr[W_i=w\mid \sigma_i,q_i]
 =\frac{(N-q_i-3)!}{N!}=\frac1{\fall N{q_i+3}}.
 \label{eq:local-batch-law}
\end{equation}
If $\mathcal T(R_i)$ is the set of ordered distinct triples avoiding $R_i$, then
\begin{equation}
 \Pr[U_i=u\mid R_i]=
 \frac{\mathbf1\{u\in\mathcal T(R_i)\}}{\fall{N-q_i}3}.
 \label{eq:local-triple-law}
\end{equation}
Conditioning also on the earlier auxiliary history and the public mask does not alter this law. Write $\operatorname{enc}(K,V)=K\Vert V$. For a 384-bit string $z$,
\begin{equation}
 \Pr[\operatorname{enc}(\sigma_{i+1})=z\mid R_i,g(a_{y(i)})]
 =\frac{\mathbf1\{z\oplus g(a_{y(i)})\in\mathcal T(R_i)\}}{\fall{N-q_i}3}.
 \label{eq:masked-state-law}
\end{equation}
For a candidate key whose unmasked halves are distinct and avoid $R_i$, summing over the remaining counter choices gives
\begin{equation}
 \Pr[K_{i+1}=k\mid R_i,g(a_{y(i)})]
 =\frac{N-q_i-2}{\fall{N-q_i}3}
 =\frac1{\fall{N-q_i}2};\label{eq:key-marginal}
\end{equation}
the probability is zero for inadmissible halves. Thus define
\begin{equation}
\begin{aligned}
 D&=(N-b)(N-b-1),\quad H=(N-b)(N-b-1)(N-b-2),\\
 \max_s\Pr[\sigma_{i+1}=s\mid R_i,g(a_{y(i)})]&\le H^{-1},\quad
 \max_k\Pr[K_{i+1}=k\mid R_i,g(a_{y(i)})]\le D^{-1}.
\end{aligned}
 \label{eq:constants}
\end{equation}
These are conditional sampling bounds in the auxiliary calculation. Lemma~\ref{lem:batch} proves their historical factorization under adaptive controls and terminal disclosure; the disclosed terminal state itself is not assigned a hidden-state bound.

\jordinote{Por qué introducimos grupos privados}{Las permutaciones privadas son un cálculo auxiliar que facilita contar distribuciones; todavía no son el generador real con tablas compartidas. Aquí obtenemos las masas del estado y de la clave tras XOR. Comprueba el denominador y la suma sobre el tercer bloque. La transferencia al generador real depende del acoplamiento posterior; no podemos saltarnos ese paso ni tratar el estado terminal revelado como secreto.}

\paragraph{Composition across requests.}
Let $X(V,\ell)=\{V+1,\ldots,V+\ell\}\pmod N$ and $X_i=X(V_i,q_i+3)$. For $1\le\ell_i,\ell_j\le L$,
\begin{equation}
 X(V_i,\ell_i)\cap X(V_j,\ell_j)\ne\varnothing
 \iff V_j-V_i\in\{1-\ell_j,\ldots,\ell_i-1\}\pmod N.
 \label{eq:interval-intersection}
\end{equation}
There are $\ell_i+\ell_j-1\le2L-1<N$ possible offsets. For $i<j$, charge the events
\[
 C^x_{ij}=\{K_i=K_j,\ X_i\cap X_j\ne\varnothing\},\qquad
 C^y_{ij}=\{K_i=K_j,\ \operatorname{set}(W_i)\cap\operatorname{set}(W_j)\ne\varnothing\}.
\]
The state-mass bound gives the first inequality below. For the second, fix the past before $W_j$: key equality is then fixed, while the fresh private list $W_j$ hits $W_i$ with probability at most $L^2/N$. Key equality has probability at most $1/D$. Therefore
\begin{equation}
 \Pr[C^x_{ij}]\le\frac{2L-1}{H},\qquad
 \Pr[C^y_{ij}]\le\frac{L^2}{ND}.
 \label{eq:main-conflicts}
\end{equation}
For adjacent groups one conditions on $\sigma_i,R_i$ before sampling $U_i$; one does not assume independence of $W_i$ and its next key. A repeated key with disjoint compatible entries remains a valid use of one permutation.

\jordinote{Dónde se esconde la dependencia entre claves}{Este es un punto delicado: la salida final de un grupo determina la clave del siguiente. Para grupos adyacentes, revisa qué se fija antes de muestrear la terna y cuándo se toma cada esperanza. No basta escribir que las claves son independientes. Tampoco cobramos toda repetición de clave como fallo: contamos posibles conflictos por intersecciones de entradas o salidas. Algunas coincidencias compatibles también se cobran conservadoramente en la cota.}

For external computations, let $\mathcal M$ be the event that some value occurs in six public requests. Sequential sampling and a union bound give
\begin{equation}
 \Pr[\mathcal M]\le N\binom r6(b/N)^6=\binom r6\frac{b^6}{N^5}.
 \label{eq:main-multiplicity}
\end{equation}
Outside $\mathcal M$, one external edge can intersect at most five public rows through its output value, plus the terminal triple. Hidden triples and input intervals contribute the other terms:
\begin{equation}
 \alpha_t=
 \underbrace{\frac6D}_{\text{disclosed values}}
 +\underbrace{\frac{3(t-1)}{D(N-b-2)}}_{\text{hidden update values}}
 +\underbrace{\frac{tL}{H}}_{\text{input intervals}}.
 \label{eq:main-alpha}
\end{equation}
The bidirectional partial-table comparison in Lemma~\ref{lem:coupling} consequently has error at most $\varepsilon_{t,r}(Q)$, where
\begin{align}
 \varepsilon_{t,r}(Q)
 &=\binom t2\left(\frac{2L-1}{H}+\frac{L^2}{ND}\right)
 +\binom r6\frac{b^6}{N^5}+Q\alpha_t\\
 &<2^{-262}+Q\,2^{-253},\qquad t\le2^{49},\ r\le2^{48}.
 \label{eq:numerics}
\end{align}
Appendix~\ref{app:composition} supplies the sequential coupling and its conditional justification, including edges computed before a future intersecting group. This comparison preserves the shared tables of the actual cipher family.

\jordinote{Qué aporta la composición completa}{Aquí se combinan los conflictos entre grupos, la multiplicidad pública y las computaciones externas. La finalidad es cubrir todo el intervalo de $2^{48}$ peticiones sin una restricción artificial a pocas solicitudes. Los verificadores comprueban las desigualdades numéricas exactas; tú debes revisar la cobertura de los eventos y el argumento de acoplamiento, que no quedan demostrados por esas cuentas automáticas.}

\paragraph{Conditional prediction.}
In one private group, fix disjoint sets of observed and target positions before sampling the group. Let $C$ be the set of values at the observed positions: $\nu$ output blocks and, when disclosed, the three Update values. Put $m=\nu+h$, $h\in\{0,3\}$. If $J$ is the set of $u$ target positions and $R_J$ their ordered values,
\begin{equation}
 \Pr[R_J=x\mid C]
 =\frac{\mathbf1\{x\text{ distinct},\ \operatorname{set}(x)\cap C=\varnothing\}}
 {\fall{N-m}u}.
 \label{eq:fixed-local-law}
\end{equation}
Here conditioning on $C$ means the complete observed values at those fixed positions. For a projection $\pi_I$ to $n$ fixed bits of one target block $B$,
\begin{equation}
 \Pr[\pi_I(B)=z\mid C]
 =\frac{2^{128-n}-|\{w\in C:\pi_I(w)=z\}|}{N-m}
 \le 2^{-n}\frac{N}{N-m}.
 \label{eq:fixed-local-bits}
\end{equation}
For example, $\nu$ known blocks give full-block baseline $1/(N-\nu)$; terminal Update disclosure changes it to $1/(N-\nu-3)$. The count is independent of the chronological direction in which a target is selected. For adaptive target selection across requests, a local count alone is insufficient; the following joint bound is used.

\jordinote{Conteo local frente a selección adaptativa}{El conteo anterior exige posiciones fijadas antes de muestrear. No debe usarse directamente cuando el observador selecciona objetivos después de ver partes de varios bloques: esa selección puede sesgar lo que queda oculto. Por eso la proposición siguiente domina la distribución conjunta de todo el intervalo. Revisa esta diferencia; ya motivó una corrección importante del argumento.}

\Needspace{6\baselineskip}
\begin{proposition}[Prediction of internal numbers]\label{prop:prediction}
Let $n$ distinct output bits be designated as targets, possibly in several requests; reveal the other output bits, processed auxiliary data and the terminal state. Throughout the task, lengths, masks, stopping and target selection depend only on the history actually disclosed. Each target bit is withheld before any observation of its value and remains hidden until the prediction. For at most $r$ requests, put
\[
 \rho_{\rm all}=\sup_{\text{allowed runs}}\prod_{i:q_i>0}\frac{N^{q_i}}{\fall{N-3}{q_i}}.
\]
This is a constant bound even for adaptive schedules. Then
\begin{equation}
 \Pr[\text{correct prediction}]\le
 \rho_{\rm all}2^{-n}+\varepsilon_{t,r}(Q)
 \le (1+2^{-54})2^{-n}+\varepsilon_{t,r}(Q).
 \label{eq:prediction}
\end{equation}

\end{proposition}
\begin{proof}
Compare the private calculation with a reference calculation in which each $R_i$ consists of $q_i$ independent uniform blocks and each $U_i$ is an independent uniform distinct triple. Keep the same state translations, public controls and terminal disclosure. For a batch with a distinct actual list, its actual/reference joint likelihood ratio is
\[
 \frac{1/\fall N{q_i+3}}{N^{-q_i}/\fall N3}
 =\frac{N^{q_i}}{\fall{N-3}{q_i}}.
\]
Otherwise the actual likelihood is zero. Multiplying these conditional likelihoods dominates the complete joint law by $\rho_{\rm all}$. Zero-output pre-update batches contribute one. Public controls are the same functions of the disclosed history in both calculations, so their indicators cancel. In the reference, every still-unrevealed output bit is independent uniform, including after disclosure of non-target bits and the independent terminal triple. The known mask cannot encode a withheld bit because controls do not inspect it. Predicting $n$ such bits succeeds with probability $2^{-n}$, even with online target selection. Density domination and the single coupling error establish the first bound. No extra historical Update triple is added to the coupling interface.

For each request the denominator divided by $N^{q_i}$ is $\prod_{j=0}^{q_i-1}(1-(3+j)/N)$. The elementary induction $\prod_j(1-x_j)\ge1-\sum_jx_j$ for $0\le x_j\le1$ gives a loss at most $\eta=b(b+3)/N<2^{-103}$. Therefore $\rho_{\rm all}\le(1-\eta)^{-r}\le(1-r\eta)^{-1}<1+2^{-54}$ for $r\le2^{48}$. The global factor is needed when target positions are selected after partially seeing blocks; a fixed-position local count would not justify that selection.
\end{proof}

Substituting (\ref{eq:numerics}) gives
\begin{equation}
 \Pr[\text{correct prediction of }n\text{ bits}]
 \le(1+2^{-54})2^{-n}+2^{-262}+Q\,2^{-253}.
 \label{eq:main-prediction}
\end{equation}

\jordinote{La pieza central y su uso en AIS}{La proposición es la pieza que conecta el cálculo con .5, .6 y la parte clásica de .7. Su prueba compara distribuciones conjuntas, incluyendo la revelación terminal, y añade una sola pérdida de acoplamiento. Revisa que los controles solo utilicen información divulgada: si una aplicación codifica el objetivo oculto en una entrada pública posterior, ella misma lo ha filtrado y se sale del experimento.}
\Needspace{5\baselineskip}
\subsection{Forward secrecy: DRG.3.5}\label{sec:secrecy}
Let the last observed block be $Z_{y,j}$. Its successor is
\begin{equation}
 Z^+_{y,j}=\begin{cases}
 \Enc_{\bar K_y}(\bar V_y+j+1),&j<q_y,\\
 \Enc_{\bar K_{y+1}}(\bar V_{y+1}+1),&j=q_y,
 \end{cases}\label{eq:successor-cases}
\end{equation}
where the second case presupposes that the next request exists. In the first case the target belongs to the same permutation sample as the observed block. For fixed observed positions, (\ref{eq:fixed-local-law})--(\ref{eq:fixed-local-bits}) give its conditional distribution. For adaptive observations, Proposition~\ref{prop:prediction} provides the joint bound. In the second case the key has undergone the transition $\phi$; Lemma~\ref{lem:coupling} is what permits the same prediction result to be applied across that boundary.

In particular, any $n$ distinct withheld bits of the successor satisfy (\ref{eq:main-prediction}). The disclosed sequence therefore provides no prediction advantage beyond the stated density and primitive-computation terms. This establishes the forward-secrecy implication of Definition~\ref{def:forward} under the declared hypotheses.

\Needspace{5\baselineskip}
\subsection{Backward secrecy: DRG.3.6}
If the first observed block is $Z_{y,j}$, its predecessor, when it exists, is
\begin{equation}
 Z^-_{y,j}=\begin{cases}
 \Enc_{\bar K_y}(\bar V_y+j-1),&j>1,\\
 \Enc_{\bar K_{y-1}}(\bar V_{y-1}+q_{y-1}),&j=1.
 \end{cases}\label{eq:predecessor-cases}
\end{equation}
Within a request the chronological relation is
\[
 Z_{y,j-1}=\Enc_{\bar K_y}\bigl(\Enc_{\bar K_y}^{-1}(Z_{y,j})-1\bigr).
\]
The unknown working key is essential; the identity does not give an inversion procedure from output alone. At the boundary the target uses the previous working key, not the disclosed output's key. It is nevertheless a withheld output covered by Proposition~\ref{prop:prediction}. Applying (\ref{eq:main-prediction}) gives the same bound as for the successor. This proves the backward-secrecy implication of Definition~\ref{def:backward}.

\Needspace{5\baselineskip}
\subsection{Enhanced backward secrecy: DRG.3.7}\label{sec:enhanced}
Disclosure of $(K^+,V^+)$ and $a$ reveals the triple in (\ref{eq:revealed-triple}). In the private calculation it excludes three values from the earlier output sample, as in (\ref{eq:fixed-local-law}) with $h=3$. The joint bound (\ref{eq:main-prediction}) already includes that disclosure, so it bounds prediction of old output itself, not merely reconstruction of a predecessor key.

For a candidate old working key $k$, write the revealed triple as $(v_1,v_2,v_3)$. Its implied final output block is
\begin{equation}
 \widehat Z_q(k)=\Enc_k(\Enc_k^{-1}(v_1)-1).
 \label{eq:candidate-last-output}
\end{equation}
The two remaining values test whether this key admits the required consecutive inverse inputs. This yields the distinct transition-inversion problem below.

The compromise boundary is essential: before the final Update, disclosure of $(\bar K,\bar V)$ permits the deterministic reconstruction
\[
 Y=\operatorname{left}_p\left(\big\Vert_{j=1}^q\Enc_{\bar K}(\bar V+j)\right).
\]
The evaluated request-level property therefore uses final-Update-before-release and erasure. It does not claim enhanced backward secrecy within an unfinished request.

\paragraph{Requirement implication.}
Equation~(\ref{eq:revealed-triple}) identifies exactly what the current state discloses about the old request. Proposition~\ref{prop:prediction} already grants that disclosure, and (\ref{eq:main-prediction}) consequently bounds recovery of the old output itself. Together with the atomicity and erasure convention in Definition~\ref{def:enhanced}, this supplies the classical enhanced-backward-secrecy argument. Transition one-wayness below is additional support, not a substitute for protecting old output. The separate infeasibility condition of AIS paragraph 214 is retained in Section~\ref{sec:discussion}.

\jordinote{Por qué .7 necesita proteger salida y no solo estado}{El BSI argumenta desde Update y el comportamiento seudoaleatorio. Aquí intentamos hacer explícito el paso que interesa: con el estado final revelado, seguir sin poder reconstruir las salidas anteriores. La unidireccionalidad de la transición no sustituye ese objetivo. Revisa el orden Update--uso y el borrado; si se compromete el estado de trabajo antes del Update, la reconstrucción mostrada sí es posible.}

\Needspace{5\baselineskip}
\subsection{Effective state and request length: DRG.3.3 and DRG.3.2}\label{sec:remaining}
Let the AES key have $x$ bits and let $r_1,r_2$ be the first two complete blocks of a request. For every candidate key $k$, the candidate working counter and consistency check are
\begin{equation}
 v_k=\Enc_k^{-1}(r_1)-1,\qquad
 \Enc_k(v_k+2)\stackrel{?}=r_2.\label{eq:key-candidate}
\end{equation}
Thus the compatible working states belong to
\[
 \mathcal S(r_1,r_2)=
 \{(k,v_k):k\in\bits^x,\ \Enc_k(v_k+2)=r_2\},
 \qquad |\mathcal S(r_1,r_2)|\le2^x.
\]
Thus the counter does not contribute a further independent factor of $2^{128}$ to a key search. AIS paragraphs 102 and 1007 identify the effective-state size with the key size for this construction. AES-128 and AES-192 are consequently below the 246-bit threshold; AES-256 provides 256 effective bits. The quantitative prediction and inversion arguments supply the computational analysis beyond this size accounting. Neither the stored-state length nor the candidate count alone establishes the other DRG.3 properties.

\jordinote{Por qué descartamos AES-128 y AES-192}{Aquí seguimos el criterio de estado efectivo del propio AIS: fijada una clave candidata, el primer bloque conocido determina el contador candidato. Revisa ese paso y las referencias AIS 102 y 1007. El conteo limita el espacio de búsqueda; por sí solo no demuestra dificultad computacional para AES-256. Esa parte necesita los argumentos posteriores.}

The request domain $I$ and the guard in Section~\ref{sec:tuples} enforce $p\le2^{19}$ bits for every successful Generate operation, establishing DRG.3.2. The reseed limit $R$ counts requests rather than bits. Neither the public counter $c$ nor its storage contributes entropy or additional effective-state bits.

\Needspace{5\baselineskip}
\subsection{Controlled additional input: DRG.3.8}
For the processed input $g(a_{y(i)})$, the law (\ref{eq:masked-state-law}) has the form
\[
 \Pr[\sigma_{i+1}=s\mid R_i,g(a_{y(i)})]
 =\Pr[U_i=\operatorname{enc}(s)\oplus g(a_{y(i)})\mid R_i,g(a_{y(i)})],
\]
where $\operatorname{enc}$ concatenates the key and counter. Thus the point-mass and key-marginal bounds in (\ref{eq:constants}) are uniform over the chosen masks. A nonempty input creates a $q_i=0$ pre-update and a $q_i>0$ output group using that same mask. With $t\le2r$, (\ref{eq:numerics}) already counts both. No independence or entropy of the input is needed.

\paragraph{Requirement implication.}
Equations (\ref{eq:masked-state-law})--(\ref{eq:key-marginal}) prove that every admitted public mask preserves the counting bounds. Lemma~\ref{lem:coupling} includes both updates, and Proposition~\ref{prop:prediction} and the inversion bounds allow public adaptive control. Therefore the established secrecy and one-wayness bounds hold uniformly over the input policies of Definition~\ref{def:additional}, including use of the same mask twice. This establishes preservation of those properties in the stated model; the XOR observation alone would not establish the composition argument.

\jordinote{Qué añadimos al argumento del XOR}{Nos inspiramos en la observación del BSI: una máscara pública no debilita una distribución uniforme por el mero XOR. Pero aquí además hay dos Update con la misma máscara y claves derivadas. Revisa que las leyes sean uniformes en la máscara elegida y que el acoplamiento cubra ambos grupos, incluso con elecciones adaptativas; el razonamiento local del XOR no basta para eso.}

\Needspace{5\baselineskip}
\subsection{Cryptographic functions and transition one-wayness: DRG.3.9}
\Needspace{6\baselineskip}
\begin{lemma}[Alignment of an inverse triple]\label{lem:inverse}
For three fixed distinct values $v_1,v_2,v_3$ and a fresh uniform permutation $\Pi$, the probability that some $x$ satisfies
\[
 \Pi(x)=v_1,\quad \Pi(x+1)=v_2,\quad \Pi(x+2)=v_3
\]
is exactly $p_*=[(N-1)(N-2)]^{-1}$.
\end{lemma}
\begin{proof}
The ordered inverse triple has $N(N-1)(N-2)$ equally likely distinct possibilities. Exactly $N$ are $(x,x+1,x+2)$, one for each $x\in\mathbb Z_N$. They are disjoint possibilities. Their ratio gives $p_*$, including wraparound.
\end{proof}
Fix the actual known $a,p$ and invert $\phi_{a,p}$ as in Definition~\ref{def:oneway}.

Suppose $(K^+,V^+)$ and the actual $a$ are known. Denote the target triple by $v_1\Vert v_2\Vert v_3=(K^+\Vert V^+)\oplus g(a)$. In an actual completed Update it is distinct. For a candidate working key $k$, let $x=\Enc_k^{-1}(v_1)$. A working-state preimage of the final Update exists exactly when
\begin{equation}
 \Enc_k^{-1}(v_2)=x+1,\qquad \Enc_k^{-1}(v_3)=x+2,\qquad
 \bar V=x-q-1.\label{eq:alignment}
\end{equation}
This equivalence, obtained by reversing (\ref{eq:post}), characterizes every working-state preimage of the final Update, not merely the original one.

\jordinote{Cualquier preimagen y la clave verdadera}{No buscamos solo la clave original: AIS pide dificultad para encontrar cualquier predecesor compatible con la entrada real. La alineación inversa caracteriza los estados de trabajo que son preimágenes del Update final. Con entrada no vacía, un estado anterior completo debe además pasar por el pre-Update; esa verificación aparece después. Revisa los índices $x,x+1,x+2$ y $\bar V=x-q-1$. La terna real proviene de alguna clave; el conteo para una etiqueta independiente no puede aplicarse sin la reducción y el control de dependencia que siguen.}

For one independent permutation label, Lemma~\ref{lem:inverse} is the exact probability that \emph{any} predecessor counter exists. An observer using $Q$ primitive computations encounters at most $Q$ labels and can return one further candidate. Appendix~\ref{app:inversion} makes this search count precise and handles adaptive stopping. It introduces the factor
\[
 \Gamma=\frac{\fall N3}{(N-b)(N-b-1)(N-b-2)}<1+2^{-110},\qquad b=4096,
\]
to account for the conditioning of the final triple on the output of its request. The resulting bounds are as follows.

Before comparing real and private success, append the three primitive evaluations that verify the proposed predecessor. This makes success an event in the same observed primitive computations in both calculations. Lemma~\ref{lem:coupling} therefore gives
\begin{equation}
 \Pr[\text{any preimage of }\phi]\le
 \varepsilon_{t,r}(Q+3)+\Gamma(Q+1)p_*\qquad(a=\varnothing).
 \label{eq:preimage}
\end{equation}
For nonempty $a$, a candidate old state first undergoes three evaluations of (\ref{eq:pre}); its candidate working key is then checked by the same three final evaluations. Allowing all labels from these six checks and one returned label gives the conservative bound
\begin{equation}
 \Pr[\text{any preimage of }\phi]\le
 \varepsilon_{t,r}(Q+6)+\Gamma(Q+7)p_* .\label{eq:preimageadditional}
\end{equation}
An adversary returning only a working predecessor was granted a larger solution set; requiring an actual old state cannot increase success. The actual public reseed counter adds a deterministic consistency check; imposing it can only restrict that solution set. The appended checks are necessary: merely comparing whether different unobserved tables admit a preimage would not be a valid coupling argument.

\paragraph{Requirement implication.}
The output $\psi$ evaluates AES and the complete $\phi$ contains its three-block Update, rather than merely an increment or XOR. For the one-wayness task in Definition~\ref{def:oneway}, equation~(\ref{eq:alignment}) characterizes all working-state candidates; Lemma~\ref{lem:inverse} counts their exact alignment probability. Appendix~\ref{app:inversion} and Lemma~\ref{lem:coupling} then give (\ref{eq:preimage}) or (\ref{eq:preimageadditional}), including verification of any returned old state. Hence any-preimage inversion has the stated classical work bound. This supplies the one-wayness component; suitability of concrete AES in this derived-key use is the cryptographic assessment discussed in Section~\ref{sec:primitive-assessment}.

\Needspace{5\baselineskip}
\subsection{Statistical inconspicuousness: DRG.3.10}
For the private output list of request $y$, sampling without replacement differs from independent uniform sampling only when the latter repeats. Fix the pre-request history $\mathcal H_y$, including its chosen length, and let $\mathcal U_q$ denote the law of $q$ independent uniform blocks. Let $Q_c$ bound the direct and inverse primitive computations used by the public input, length and stopping controller over the interval. The test itself observes only output. Then
\begin{align}
 \Delta(\mathcal L(R_y\mid\mathcal H_y),\mathcal U_{q_y})
 &\le\binom{q_y}{2}/N\le\binom b2/N,\\
 \Delta(\text{public output},\text{ideal output})
 &\le r\binom b2/N+\varepsilon_{t,r}(Q_c)<2^{-57}+Q_c\,2^{-253}.
 \label{eq:statistics}
\end{align}
The uniform per-request bound also covers adaptive lengths under the same public request rule. Truncation cannot increase distance. By Definition~\ref{def:statistics}, equation~(\ref{eq:statistics}) bounds the change in acceptance probability of every output-only test by $2^{-57}+Q_c2^{-253}$ over the entire interval. In particular the bound is $2^{-57}$ for a public controller making no primitive computations ($Q_c=0$). Being an output-only test does not by itself impose that restriction on the controller. This is the model-based statistical-inconspicuousness argument for DRG.3.10. It is separate from the computational prediction bounds and does not replace checks of an implementation.

\jordinote{La precisión sobre el controlador en .10}{No confundas un test que solo mira salida con un controlador que no calcula AES. Si el controlador hace $Q_c$ computaciones, ese presupuesto aparece en el error de comparación; el número $2^{-57}$ aislado exige $Q_c=0$. Revisa también que la cota estadística de un intervalo no se presente como fuerza de predicción de 57 bits. Esta distinción corrige una sobreafirmación anterior.}

\Needspace{5\baselineskip}
\subsection{Source premise and initialized entropy: DRG.3.1 and DRG.3.4}\label{sec:entropy}
The source assumption in Section~\ref{sec:model} supplies the distribution of the entropy input. We now verify its algorithmic propagation through the two seeding maps, both of which are XOR translations with a mask fixed by the known history. For a finite random variable $X$, use
\[
 H_{\min}(X)=-\log_2\max_x\Pr[X=x],\qquad
 H(X)=-\sum_x\Pr[X=x]\log_2\Pr[X=x].
\]
\Needspace{6\baselineskip}
\begin{lemma}[Translation and projection]\label{lem:entropy}
For fixed $m\in\bits^{384}$ and $W=e\oplus m=K\Vert V$,
\begin{align}
 \hm(W)&=\hm(e),& H(W)&=H(e),\label{eq:translate}\\
 \hm(K)&\ge\hm(e)-128,& H(K)&\ge H(e)-128.\label{eq:project}
\end{align}
The inequalities also hold separately at every fixed history where the mask is determined.
\end{lemma}
\begin{proof}
XOR by $m$ is its own inverse, so it permutes the point masses. Each key has exactly $2^{128}$ counter extensions; consequently $\Pr[K=k]\le2^{128}\max_w\Pr[W=w]$. Taking negative binary logarithms gives the min-entropy bound. The Shannon chain rule gives $H(W)=H(K)+H(V\mid K)\le H(K)+128$. First fixing the history proves the conditional statements without assuming marginal independence.
\end{proof}

For the two maps the masks are $m=\pad(z)\oplus T(0,0)$ and $m=\pad(z)\oplus T(K,V)$. Under conditional uniform input,
\begin{equation}
 \Pr[K=k,V=v\mid\text{history}]=2^{-384},\qquad
 \Pr[K=k\mid\text{history}]=2^{128}2^{-384}=2^{-256}.
 \label{eq:uniform-initial-state}
\end{equation}
Thus both initialized-key entropies equal 256 bits. For DRG.3.4, Lemma~\ref{lem:entropy} and (\ref{eq:uniform-initial-state}) give $256\ge240$ min-entropy bits in the initial effective key under the assumed source premise, including after reseeding a known old state. DRG.3.1, admissibility of the actual source, is assumed rather than inferred from this algebra. The general projection bounds give initial key min-entropy 240 from seed min-entropy 368, or Shannon entropy 250 from seed Shannon entropy 378 with the required stationarity. Those sufficient initial-entropy bounds do not extend the main composition result beyond its uniform 384-bit premise.

For comparison, a uniform key on $2^{122}$ values and an independent uniform counter give $H_{\min}(K,V)=250$ but $H_{\min}(K)=122$. Pulling this distribution back through the XOR translation preserves its entropy. This refutes an inference from 250 joint seed bits alone to sufficient effective-key entropy; it is not a counterexample to the stated initialization premise.

\jordinote{Entropía efectiva y alcance de la premisa}{La fuente suficiente se asume, pero sí hay que justificar qué hace el algoritmo con ella. Para uniform384, XOR da un estado uniforme y una clave uniforme256. El ejemplo de 250 bits conjuntos muestra por qué no basta contar entropía almacenada. Revisa que los umbrales generales 368/378 solo se usen para entropía inicial, no para extender silenciosamente toda la prueba de composición.}

\section{Conclusions and open questions}\label{sec:discussion}
For the stated AES-256 profile, the describing tuples lead to a requirement-by-requirement algorithmic argument. The request bound and effective-state size follow from the specification and state accounting. The shared-key composition estimate and prediction bound establish the classical secrecy implications, including protection of earlier output after terminal-state disclosure. The same input processing preserves these estimates under public adaptive additional input. Inverse alignment gives a bound for finding any predecessor of the complete request transition; output sampling gives the statistical-distance estimate. Table~\ref{tab:traceability} records the results used for each implication.

\jordinote{Cómo revisar el cierre requisito por requisito}{Comprueba esta tabla en ambos sentidos: cada requisito debe tener una definición y una justificación; cada resultado utilizado debe tener una hipótesis identificable. La comparación de rigor con Hash\_DRBG está en el registro de revisión, para no convertir el artículo en una discusión sobre cómo se ha redactado. Una conclusión condicional no puede perder sus condiciones al resumirse.}

\begingroup\small\setlength{\tabcolsep}{5pt}
\begin{longtable}{@{}L{17mm}L{126mm}@{}}
\caption{Traceability of the DRG.3 argument under the stated hypotheses.}\label{tab:traceability}\\
\toprule DRG.3 & Mathematical justification\\\midrule\endfirsthead
\toprule DRG.3 & Mathematical justification\\\midrule\endhead
\bottomrule\endfoot
.1 & Admissibility of the selected TRNG source is assumed.\\[2pt]
.2 & The request domain and Generate guard enforce $p\le2^{19}$.\\[2pt]
.3 & Candidate-counter reconstruction (\ref{eq:key-candidate}) and the AIS effective-state accounting give 256 bits.\\[2pt]
.4 & Lemma~\ref{lem:entropy} and (\ref{eq:uniform-initial-state}) give 256 initialized-key entropy bits under conditional uniform384 input.\\[2pt]
.5 & Successor identity (\ref{eq:successor-cases}), Proposition~\ref{prop:prediction} and Lemma~\ref{lem:coupling} give (\ref{eq:main-prediction}).\\[2pt]
.6 & Predecessor identity (\ref{eq:predecessor-cases}) and the same prediction result give the corresponding backward bound.\\[2pt]
.7 & Terminal-state identity (\ref{eq:revealed-triple}) and Proposition~\ref{prop:prediction} bound classical prediction of earlier output at the specified atomicity and erasure boundary.\\[2pt]
.8 & Equations (\ref{eq:masked-state-law})--(\ref{eq:key-marginal}) and Lemma~\ref{lem:coupling} preserve the prediction and inversion bounds for the admitted input policies.\\[2pt]
.9 & Lemma~\ref{lem:inverse}, (\ref{eq:alignment}) and Appendix~\ref{app:inversion} give (\ref{eq:preimage})--(\ref{eq:preimageadditional}); concrete-primitive suitability remains a cryptographic assessment.\\[2pt]
.10 & Equation~(\ref{eq:statistics}) bounds output-only statistical tests by $2^{-57}+Q_c2^{-253}$.\\
\end{longtable}
\endgroup

The mathematical conclusions hold in the declared permutation model. Section~\ref{sec:primitive-assessment} explains the cryptographic hypothesis used when applying that model to concrete AES-256, distinguishing single-key evidence, related-key access and generic quantum search. The resulting conformity argument is conditional on that assessment, as well as on source sufficiency and the selected compromise boundary. The evidence does not establish a universal quantum lower bound or a reduction from every attack on CTR\_DRBG to single-key AES key recovery. Such stronger claims would require further work; they are not consequences of the classical inequalities.

\jordinote{Qué significa lo que sigue pendiente}{El ejemplo Hash\_DRBG usa hipótesis criptográficas justificadas, incluidas algunas no estándar, y no aporta una prueba cuántica separada. Por eso no debemos exigir aquí un teorema de ese tipo como nuevo requisito para estar al mismo nivel. Sí debemos distinguir las cotas clásicas demostradas de la valoración criptográfica adicional que requiere AIS 214. La revisión comparada identifica exactamente esa diferencia.}

\paragraph{Extension to the derivation-function configuration.}
With DF, Instantiate and Reseed instead supply Update with
\[
 \DF(e\Vert\mathsf{nonce}\Vert z,384)
 \quad\text{and}\quad \DF(e\Vert z,384),
\]
respectively, under the nonce and alignment rules of SP 800-90A. A known processed additional-input mask can be used in the generation argument. The different initialization map still needs an entropy justification.

Let $E$ be seed input, $Z$ the disclosed history including the old state at reseeding, and $F_z$ the seeding map followed by key projection. For each history $z$ of positive probability and largest fibre $M_z$ on the conditional support,
\begin{equation}
 H_{\min}(F_z(E)\mid Z=z)\ge H_{\min}(E\mid Z=z)-\log_2M_z.\label{eq:fibre}
\end{equation}
Indeed, each output mass sums at most $M_z$ input masses. A useful fibre bound or another justified effective-key entropy estimate would establish this initialization implication. Applying the generation bounds would additionally require an argument for the resulting joint distribution of the initial key and counter. SP 800-90B table 1 lists AES-256 \DF\ with 256-bit output, whereas CTR\_DRBG requests 384. The requested length enters the BCC input, so $\operatorname{left}_{256}\DF(e,384)$ is not identically $\DF(e,256)$. The table alone does not establish the required implication.

\jordinote{La diferencia con el argumento DF del BSI}{El estudio estima entropía mediante un modelo de aplicación aleatoria para DF. Nosotros no hemos desarrollado aquí una justificación de entropía efectiva para ese mapa y configuración. Revisa la diferencia entre pedir 256 y 384 bits a Block\_Cipher\_df y el alcance real de la tabla de la 90B. No estamos afirmando que DF falle: estamos separando la implicación que aún necesitaría justificación.}

Each result concerns one finite reseed interval; bounds for finitely many intervals may be summed after re-establishing the fresh source premise. The separate extension to more general seed distributions requires a composition argument for those distributions. Side channels, faults and state disclosure during an unfinished request are not part of the algorithmic statement proved here.

\appendix
\section{Composition of requests under shared key labels}\label{app:composition}
This appendix supplies the conditional distributions and partial-table comparison used in (\ref{eq:numerics}). A pre-update is a group with no public output; an output request and its final update form a group with $q_i$ public blocks. A public XOR mask is applied after the three update values. There are at most two groups per request, including when the same mask is used twice.

Use the constants in (\ref{eq:constants}).
Write $\fall{x}{j}=x(x-1)\cdots(x-j+1)$. A batch consists of $q_i$ visible output blocks followed by the three Update blocks; $0\le q_i\le b$. Nonempty additional input creates a pre-update batch with $q_i=0$, then an output batch with $q_i\ge1$. Each batch applies a fixed public XOR mask $g(a_{y(i)})$ to its Update triple to obtain the next state. The same mask in the two batches of one request is retained. There are $t\le2R$ batches and $r\le R$ output requests.

For an auxiliary calculation, give batch $i$ a private independent uniform permutation $\Pi_i$, even if its key label equals a previous label. Its consecutive inputs form $X_i$; its output list is $W_i=R_i\Vert U_i$, where $|R_i|=q_i$ and $|U_i|=3$. If $\sigma_i=(K_i,V_i)$, then
\begin{equation}
 \sigma_{i+1}=\operatorname{split}_{256,128}(U_i\oplus g(a_{y(i)})).\label{eq:batchstate}
\end{equation}
The private permutations define an auxiliary probability calculation. The following comparison relates it to the shared-key permutation family.

\Needspace{6\baselineskip}
\begin{lemma}[Conditional batch law]\label{lem:batch}
Conditional on the past and the selected length, $W_i$ is uniform among distinct ordered lists of length $q_i+3$. Conditional on its public $R_i$, $U_i$ is uniform among distinct ordered triples avoiding $R_i$. Hence any specified historical state has conditional mass at most $1/H$, and any specified historical key has mass at most $1/D$.
\end{lemma}
\begin{proof}
The inputs are distinct because $q_i+3<N$. Exactly $(N-q_i-3)!$ of the $N!$ permutations produce any specified distinct list at these inputs. Its probability is therefore $1/\fall{N}{q_i+3}$, independently of the input interval and incoming state. Fixing $R_i$ leaves $\fall{N-q_i}{3}$ equiprobable triples. A specified pair of first two triple values has probability at most $1/\fall{N-q_i}{2}$; translation by the public mask changes no mass. Both denominators are at least $H$ and $D$ respectively.

More explicitly, for fixed public choices and lists the joint likelihood factors into the factors $1/\fall{N}{q_i+3}$. Thus, conditional on all public lists and choices, the hidden $U_i$ factor as independent uniform triples on their respective allowed sets. Adaptive choices do not add a hidden restriction: each choice was a function of already public observations and the observer's independent external tables. Revealing the terminal state fixes $U_t$ alone through (\ref{eq:batchstate}). Historical states still have the stated mass bounds. The initial $\sigma_1$ is uniform and independent of the primitive by Lemma~\ref{lem:entropy} and the uniform source premise.
\end{proof}

\jordinote{Punto crítico del apéndice: condicionamientos}{Aquí conviene revisar la factorización con calma: qué variables quedan ocultas después de fijar transcripción, controles, tablas externas y estado terminal. La última terna está determinada por el compromiso y debe tratarse aparte. El argumento necesita que las decisiones públicas dependan solo del historial observado; no basta invocar independencia de permutaciones sin comprobar qué condicionamientos la conservan.}

\Needspace{5\baselineskip}
\subsection{Repeated evaluations, not repeated labels}
\Needspace{6\baselineskip}
\begin{lemma}[A full-interval partial-table bound]\label{lem:coupling}
For $t$ batches, $r$ output requests and $Q$ external direct or inverse primitive computations, define
\begin{align}
 \delta_{\rm priv}&=\binom t2\left(\frac{2L-1}{H}+\frac{L^2}{ND}\right),\label{eq:priv}\\
 \delta_{\rm mult}&=\binom r6\frac{b^6}{N^5},\label{eq:mult}\\
 \alpha_t&=\frac6D+\frac{3(t-1)}{D(N-b-2)}+\frac{tL}{H},\label{eq:alpha}\\
 \varepsilon_{t,r}(Q)&=\delta_{\rm priv}+\delta_{\rm mult}+Q\alpha_t.\label{eq:epsilon}
\end{align}
The strengthened observation of the real ideal-cipher generator and of the private-batch calculation can be coupled to agree except with probability at most $\varepsilon_{t,r}(Q)$.
\end{lemma}
\begin{proof}
For $i<j$, charge a private conflict when $K_i=K_j$ and either $X_i\cap X_j\ne\varnothing$ or the output sets of $W_i,W_j$ intersect. An input overlap requires $K_j=K_i$ and at most $2L-1$ possible counter values, since $V_j-V_i$ must be one of the offsets from $-L+1$ to $L-1$. Lemma~\ref{lem:batch} gives probability at most $(2L-1)/H$. For adjacent batches, condition on $\sigma_i,R_i$ and then count the allowed $U_i$; do not assume $\sigma_{i+1}$ independent of $W_i$.

For an output overlap, condition just before $W_j$ is generated. The equality of keys and $W_i$ are already fixed, while $W_j$ is a fresh uniform list. The conditional chance of an intersecting value is at most $L^2/N$. The equality of keys has probability at most $1/D$: when adjacent, condition on $\sigma_i,R_i$ before sampling the translated $U_i$. This proves $L^2/(ND)$ per pair and (\ref{eq:priv}). Public masks do not change either count, since translation is a bijection on the possible next states.

Let $\mathcal M$ be the event that some 128-bit value appears in six distinct public output requests. For a fixed value and six indices the sequential conditional probability is at most $(b/N)^6$, even for adaptive lengths chosen before the corresponding list. Union over values and indices gives (\ref{eq:mult}). Outside $\mathcal M$, any value lies in at most five public output rows.

Now expose one external table edge $(k,x,y)$, whether learned by a direct or inverse computation. In the private calculation the external tables are independent of the private hidden triples, conditional on the public lists, choices and terminal state. A conflict of inputs requires, for some historical row, $K_i=k$ and one of at most $L$ values of $V_i$; its probability is at most $tL/H$. For outputs, $y$ lies in at most five public rows, each needing $K_i=k$, giving $5/D$. The terminal Update triple is known and contributes at most one further row, $1/D$. In each other hidden triple, the row key and its own $U_i$ are independent conditional on the public data, so
\[
 \Pr[K_i=k,\ y\in U_i]\le\frac{3}{D(N-b-2)}.
\]
The exceptional last triple is handled separately, not treated as hidden. Summing gives (\ref{eq:alpha}). These conditional counts also cover an external edge computed before its conflicting batch: fix the entire public transcript and external tables first, then count the remaining hidden triples. Adaptive stopping uses at most the stated number of edges and rows.

For completeness, compatibility of final tables alone is not sufficient for a coupling. Maintain a bidirectional partial table for each private batch and for the external cipher, and a merged partial table per key for the real cipher. For a new input, a hit in another record under that key is charged as a conflict. Otherwise the private candidate output is uniform among values unused in its own record. If it is used by another record under the same key, charge a conflict and sample the real value uniformly from the still admissible values of the merged table. Each unused merged output has identical probability under this construction; thus the real answer has exactly the conditional law of a uniform permutation. Inverse queries use the symmetric construction. Queries repeated inside the same record are deterministic. This shared construction is used only up to the first conflict; thereafter continue each calculation separately with its own correct conditional law. Before that conflict both observations agree. The pair and edge counts above bound all first-conflict possibilities, establishing (\ref{eq:epsilon}).
\end{proof}

\jordinote{El acoplamiento debe ser secuencial}{Que las tablas finales sean compatibles no demuestra por sí solo que las dos ejecuciones tengan las distribuciones correctas. Revisa el muestreo de cada nueva arista, las consultas inversas, los aciertos en tablas anteriores y qué ocurre al primer conflicto. Antes de ese evento deben coincidir las observaciones y cada marginal debe seguir siendo la de una permutación uniforme.}

At $t\le2^{49}$ groups and $r\le2^{48}$ output requests, the exact expressions above satisfy
\[
 \delta_{\rm priv}+\delta_{\rm mult}<2^{-262},\qquad \alpha_t<2^{-253}.
\]
Substitution into (\ref{eq:epsilon}) gives (\ref{eq:numerics}). The bounds are monotone in the permitted numbers of rows and primitive computations, so they also cover earlier stopping.

\section{Adaptive stopping and transition inversion}\label{app:inversion}
There is a small conditioning issue if the observer stops generation adaptively after seeing an output. The terminal triple then avoids a publicly selected output set and need not be marginally independent of the external cipher. Let $G$ denote the full public generation transcript before disclosure of that triple and $O$ the external cipher. Lemma~\ref{lem:batch} gives
\[
 \Pr[U_t=u\mid O,G]\le\frac1H.
\]
Compare with the same $(O,G)$ and a distinct uniform triple $U_*$ independent of them. The joint density ratio is at most
\begin{equation}
 \Gamma=\frac{\fall N3}{H}.\label{eq:gamma}
\end{equation}
Thus any success probability in the original private calculation is at most $\Gamma$ times that in the independent-target calculation. The transcript $G$ can be simulated with public independent batch samples and the observer's own primitive computations; it gives no free extra primitive tables. Giving $U_*$ to the observer before those computations only strengthens it.

For absent additional input, an observer using $Q$ primitive computations touches at most $Q$ external key labels and can return one additional label. In the independent-target calculation, even granting the \emph{entire} permutation at each touched label, each newly selected label satisfies the alignment with probability $p_*$. Adaptation between independent labels does not change that probability. Union gives $(Q+1)p_*$. This table grant is used only for this search count, not in the partial-table coupling.

The full-table grant in this search argument concerns only the private external cipher. The real/private comparison still counts individual direct or inverse edges. With no pre-update, verification of a returned predecessor adds three primitive evaluations. With a pre-update, it adds three evaluations to compute the working state and three to check its final update. Applying the comparison to this verified event gives (\ref{eq:preimage}) and (\ref{eq:preimageadditional}); it is not sufficient to compare whether two unobserved tables admit a predecessor.

\jordinote{Revisión de la reducción de inversión}{En este apéndice, $\Gamma$ corrige el sesgo de parar tras observar salida. Revisa por qué es un factor multiplicativo. La concesión de tablas completas solo se usa en el conteo privado de etiquetas, no en el acoplamiento real. Las tres o seis evaluaciones añadidas verifican el candidato devuelto: permiten comparar el mismo evento observable, no una existencia abstracta de preimágenes.}

\clearpage
\begingroup\small
\bibliographystyle{plainurl}
\bibliography{references}
\endgroup
\end{document}
